\chapter{संख्या के रूप में भिन्नें}\label{ch:g5:fractions}

\omterm{def:g4:fractions:def}{भिन्न} सिर्फ़ किसी पाई का चित्र नहीं है: वह सचमुच की एक
\emph{संख्या} है, जिसका संख्या रेखा पर अपना स्थान है --- कभी-कभी
ठीक वही स्थान जो किसी दशमलव संख्या का है। यह अध्याय दोनों दुनियाएँ
जोड़ता है। (भिन्नों की गणित के पूरे नियम \cref{ch:g6:fractions}
और उसके आगे आते हैं।)

\section{रेखा पर भिन्नें, फिर से}

\begin{example}[रखना और पढ़ना]\label{ex:g5:fractions:place}
पाँचवें भागों में कटी रेखा पर $\frac{7}{5}$, $0$ से $7$ कदम पर
है: $1$ के आगे (जो $\frac55$ है), $1 + \frac25$ पर। हर \omterm{def:g4:fractions:def}{भिन्न}
रेखा पर कहीं न कहीं रहती है; हर वही हो तो बड़ा \omterm{def:g4:fractions:def}{अंश} यानी और दाईं ओर।
\end{example}

\begin{proposition}[एक ही बिंदु, कई नाम]\label{prop:g5:fractions:equal}
हर हिस्से को दो में (या तीन में\,\dots) काट देने से बिंदु नहीं
हिलता:
\[
\frac{1}{2} = \frac{2}{4} = \frac{3}{6} = \frac{5}{10},
\qquad
\frac{2}{3} = \frac{4}{6} = \frac{20}{30}.
\]
\omterm{def:g4:fractions:def}{अंश} और हर दोनों को एक ही संख्या से गुणा (या भाग) करने पर उसी
\omterm{def:g4:fractions:def}{भिन्न} का दूसरा नाम मिल जाता है।
\end{proposition}
\admitted

\begin{omfigure}
\begin{tikzpicture}[scale=0.62]
  \foreach \i in {0,1} \draw[omExo] (\i*4,0) rectangle (\i*4+4,1.0);
  \fill[omDef!40] (0.06,0.06) rectangle (3.94,0.94);
  \draw[thick] (0,0) rectangle (8,1.0);
  \node[right, font=\small] at (8.3,0.5) {$\frac12$};
  \begin{scope}[yshift=-1.7cm]
  \foreach \i in {0,...,3} \draw[omExo] (\i*2,0) rectangle (\i*2+2,1.0);
  \fill[omDef!40] (0.06,0.06) rectangle (1.94,0.94);
  \fill[omDef!40] (2.06,0.06) rectangle (3.94,0.94);
  \draw[thick] (0,0) rectangle (8,1.0);
  \node[right, font=\small] at (8.3,0.5) {$\frac24$};
  \end{scope}
  \begin{scope}[yshift=-3.4cm]
  \foreach \i in {0,...,9} \draw[omExo] (\i*0.8,0) rectangle
    (\i*0.8+0.8,1.0);
  \foreach \i in {0,...,4} \fill[omDef!40] (\i*0.8+0.05,0.06)
    rectangle (\i*0.8+0.75,0.94);
  \draw[thick] (0,0) rectangle (8,1.0);
  \node[right, font=\small] at (8.3,0.5) {$\frac{5}{10}$};
  \end{scope}
\end{tikzpicture}

{\small तीन पट्टियाँ, तीन नाम, एक ही मात्रा: $\frac12 = \frac24 =
\frac{5}{10}$।}
\end{omfigure}

\section{भिन्नें और दशमलव संख्याएँ}

\begin{proposition}[दशमलव नाम वाली भिन्नें]\label{prop:g5:fractions:decimal}
दसवें, सौवें और हज़ारवें भागों की भिन्नें दशमलव संख्याएँ
\emph{हैं ही}:
\[
\frac{3}{10} = 0.3, \qquad \frac{27}{100} = 0.27 .
\]
कुछ और भिन्नों के भी दशमलव नाम होते हैं, जो दसवें या सौवें भागों
में बदलकर मिल जाते हैं:
\[
\frac{1}{2} = \frac{5}{10} = 0.5, \qquad
\frac{1}{4} = \frac{25}{100} = 0.25, \qquad
\frac{3}{4} = 0.75, \qquad
\frac{1}{5} = \frac{2}{10} = 0.2 .
\]
पर सबके नहीं: $\frac13 = 0.333\dots$ कभी ख़त्म नहीं होता ---
\omterm{def:g4:fractions:def}{भिन्न} ही उसका इकलौता ठीक नाम है।
\end{proposition}
\admitted

\begin{omfigure}
\begin{tikzpicture}[scale=1.55]
  \draw[-Stealth] (-0.2,0) -- (7.6,0);
  \foreach \i in {0,...,28} \draw ({\i*0.25},-0.05) -- ({\i*0.25},0.05);
  \foreach \i in {0,4,...,28} \draw ({\i*0.25},-0.09) -- ({\i*0.25},0.09);
  \foreach \x in {0,...,7} \node[below, font=\small] at (\x,-0.1) {$\x$};
  \fill[omDef] (0.5,0) circle (1.6pt)
    node[above, font=\small] {$\frac12 = 0.5$};
  \fill[omThm] (2.75,0) circle (1.6pt)
    node[above, font=\small] {$\frac{11}{4} = 2.75$};
  \fill[omMeth] (6.25,0) circle (1.6pt)
    node[above, font=\small] {$\frac{25}{4} = 6.25$};
\end{tikzpicture}

{\small एक रेखा, दो भाषाएँ: हर बिंदु का एक भिन्न-नाम और एक
दशमलव नाम हो सकता है।}
\end{omfigure}

\begin{example}[जो नाम आसान पड़े, वही चुनो]\label{ex:g5:fractions:choose}
$12$ के तीन \omterm{def:g3:division:half}{चौथाई} देने हों: \omterm{def:g4:fractions:def}{भिन्न} वाला नाम आसान है, $12$ का
$\frac34$, $9$ है। $\frac34$ और $0.8$ की तुलना करनी हो: दशमलव नाम
आसान है, $0.75 < 0.8$। दोनों भाषाएँ आने का यही फ़ायदा है।
\end{example}

\section{भिन्नों की तुलना}

\begin{method}[तुलना के तीन औज़ार]\label{met:g5:fractions:compare}
\begin{enumerate}
  \item \emph{हर एक जैसा}: जिसमें ज़्यादा हिस्से, वही बड़ी,
        $\frac57 > \frac37$;
  \item \emph{\omterm{def:g4:fractions:def}{अंश} एक जैसा}: जिसके हिस्से बड़े, वही बड़ी, यानी
        \emph{छोटा} हर जीतता है: $\frac35 > \frac38$ (पाँचवें
        भाग आठवें भागों से बड़े होते हैं);
  \item \emph{ठिकाने}: हर एक की तुलना $\frac12$ या $1$ से करो:
        $\frac38 < \frac12 < \frac59$, इसलिए $\frac38 < \frac59$।
\end{enumerate}
\end{method}

\begin{example}\label{ex:g5:fractions:compare}
पिज़्ज़ा किसने ज़्यादा खाया: अली ने $\frac58$ या बिया ने
$\frac54$? बिया की \omterm{def:g4:fractions:def}{भिन्न} $1$ से बड़ी है (पाँच \omterm{def:g3:division:half}{चौथाई} एक पूरे
पिज़्ज़े से ज़्यादा हैं), अली की $1$ से छोटी: बिया ने ज़्यादा
खाया --- असल में एक पूरे पिज़्ज़े से भी ज़्यादा।
\end{example}

\section{अभ्यास}

\begin{exercise}[$\star$]\label{exo:g5:fractions:1}
$0$ से $2$ तक पाँचवें भागों में कटी संख्या रेखा बनाओ, और रखो:
$\frac25$; \ $\frac55$; \ $\frac85$; \ $\frac{10}{5}$।
\end{exercise}

\begin{exercise}[$\star$]\label{exo:g5:fractions:2}
बराबर नाम पूरे करो:
$\frac12 = \frac{?}{8}$; \quad
$\frac34 = \frac{?}{8}$; \quad
$\frac{2}{5} = \frac{4}{?}$; \quad
$\frac{6}{9} = \frac{2}{?}$।
\end{exercise}

\begin{exercise}[$\star$]\label{exo:g5:fractions:3}
कौन-सी भिन्नें पूर्ण संख्याओं के नाम हैं? जहाँ हों, वह पूर्ण
संख्या बताओ: $\frac82$; \quad $\frac{9}{4}$; \quad
$\frac{15}{3}$; \quad $\frac{20}{5}$।
\end{exercise}

\begin{exercise}[$\star$]\label{exo:g5:fractions:4}
दशमलव के रूप में लिखो: $\frac{7}{10}$; \quad $\frac{41}{100}$; \quad
$\frac12$; \quad $\frac34$; \quad $\frac15$।
\end{exercise}

\begin{exercise}[$\star$]\label{exo:g5:fractions:5}
\omterm{def:g4:fractions:def}{भिन्न} के रूप में लिखो (दसवें या सौवें भाग): $0.9$; \quad $0.13$;
\quad $2.5$; \quad $0.75$ (फिर उसका \omterm{def:g3:division:half}{चौथाई} वाला नाम बताओ)।
\end{exercise}

\begin{exercise}[$\star$]\label{exo:g5:fractions:6}
तुलना करो, और बताओ कौन-सा औज़ार लगाया
(\cref{met:g5:fractions:compare}):
\[
\frac47 \;?\; \frac67, \qquad
\frac35 \;?\; \frac37, \qquad
\frac25 \;?\; \frac{7}{12}, \qquad
\frac98 \;?\; \frac{11}{12} .
\]
\end{exercise}

\begin{exercise}[$\star$]\label{exo:g5:fractions:7}
दशमलव नामों से क्रम में लगाओ: $\frac12$; \ $0.4$; \ $\frac34$; \
$0.6$; \ $\frac15$।
\end{exercise}

\begin{exercise}[$\star$]\label{exo:g5:fractions:8}
$28$ विद्यार्थियों की कक्षा में एक \omterm{def:g3:division:half}{चौथाई} कोई वाद्य बजाते हैं, और
आधे कोई खेल खेलते हैं। हर समूह में कितने विद्यार्थी हैं? कौन-सा समूह बड़ा है?
\end{exercise}

\begin{exercise}[$\star$]\label{exo:g5:fractions:9}
हर \omterm{def:g4:fractions:def}{भिन्न} को एक समय से मिलाओ: $\frac14$ घंटा, $\frac12$ घंटा,
$\frac34$ घंटा, $\frac{1}{60}$ घंटा --- और $30$ मिनट, $45$ मिनट,
$1$ मिनट, $15$ मिनट।
\end{exercise}

\begin{exercise}[$\star\star$]\label{exo:g5:fractions:10}
टॉम ने अपनी $50$ cL की बोतल का $\frac35$ पिया; अना ने अपनी
$60$ cL की बोतल का $\frac12$। किसने ज़्यादा पिया? (दोनों मात्राएँ
cL में निकालो --- यहाँ केवल भिन्नों की तुलना काफ़ी नहीं, और बताओ क्यों।)
\end{exercise}

\begin{exercise}[$\star\star$]\label{exo:g5:fractions:11}
$\frac12$ और $1$ के ठीक बीच की कोई \omterm{def:g4:fractions:def}{भिन्न} बताओ; फिर $\frac12$ और
$\frac34$ के ठीक बीच की एक। (दशमलव नाम मदद कर सकते
हैं।)
\end{exercise}
